\documentclass[../../../include/open-logic-section]{subfiles}

\begin{document}
	
\olfileid[id]{sth}{z}{nat}
\olsection{Memilih Bilangan Asli Kita}

%Let us take it for granted that the informal principle \stagesinf{}
%can be justified. It is also, though, worth commenting on the
%particular axiom we extracted from that informal principle. 

Dalam \olref[infinity-again]{defnomega}, kita secara eksplisit mendefinisikan
ungkapan ``bilangan asli.'' Bagaimana Anda seharusnya memahami ketetapan ini?
Ketetapan itu bukan klaim metafisis, melainkan sekadar keputusan untuk
\emph{memperlakukan} himpunan tertentu sebagai bilangan asli. Kita telah
menyinggung alasan untuk berpikir demikian dalam
\olref[sfr][rel][ref]{sec}, \olref[sfr][arith][ref]{sec}, dan
\olref[sfr][infinite][dedekindsproof]{sec}. Namun, kita dapat mengemukakan
alasan-alasan itu dengan lebih tajam.

Aksioma Ketakhinggaan kita mengikuti \citet{VonNeumann1925}. Akan tetapi,
berikut aksioma lain yang dapat kita pilih sebagai gantinya:

\begin{defish}
\emph{Aksioma Ketakhinggaan Zermelo
\citeyear{Zermelo1908Untersuchungen}.} Terdapat himpunan~$A$ sedemikian
sehingga $\emptyset\in A$ dan $(\forall x\in A)\{x\}\in A$.
\end{defish}

Seandainya kita menggunakan aksioma Zermelo alih-alih Aksioma Ketakhinggaan
yang diilhami von Neumann, kita juga akan memperoleh suatu himpunan yang tak
hingga Dedekind dan, dengan demikian, suatu aljabar Dedekind. Dalam pendekatan
Zermelo, anggota istimewa aljabar kita tetap $\emptyset$ (pengganti~$0$),
tetapi injeksinya diberikan oleh pemetaan $x\mapsto\{x\}$, bukan
$x\mapsto x\cup\{x\}$. Konsekuensi paling sederhana ialah bahwa Zermelo
memperlakukan $2$ sebagai $\{\{\emptyset\}\}$, sedangkan kita (mengikuti von
Neumann) memperlakukan $2$ sebagai
$\{\emptyset,\{\emptyset\}\}$.

Mengapa memilih satu aksioma Ketakhinggaan dan bukan yang lain? Alasan praktis
utama ialah bahwa pendekatan von Neumann dapat ``diperluas'' dengan baik untuk
menangani bilangan transfinit. Kita akan menjelajahinya mulai
\olref[ordinals][]{chap}. Namun, dari sudut pandang sederhana
\emph{mengerjakan aritmetika}, kedua pendekatan sama-sama memadai. Jadi, jika
seseorang mengatakan bahwa bilangan asli \emph{adalah} himpunan, pertanyaan
yang jelas ialah: \emph{Himpunan yang mana?}

Pertanyaan persis ini menjadi terkenal berkat \citet{Benacerraf1965}. Namun,
patut ditekankan bahwa pertanyaan itu hanyalah contoh paling terkenal dari
suatu fenomena yang telah kita jumpai berkali-kali. Pokok dasarnya adalah
sebagai berikut. Teori himpunan memberi kita cara untuk \emph{menyimulasikan}
berbagai jenis entitas ``intuitif'': bilangan real, rasional, bulat, dan asli;
tetapi juga pasangan terurut, fungsi, dan relasi. Akan tetapi, teori himpunan
tidak pernah memberi kita pilihan simulasi yang \emph{unik}. \emph{Selalu}
ada pilihan alternatif yang---secara langsung---dapat melayani kita dengan
sama baiknya.

\end{document}
